% ============================================================
% Chapitre 11 : Methodes Non-Parametriques
% ============================================================
\chapter{M\'ethodes Non-Param\'etriques}
\label{ch:nonparam}

\begin{center}
\textit{<<~Ne faites pas d'hypoth\`eses inutiles.~>>} --- Guillaume d'Ockham
\end{center}

% ------------------------------------------------------------
\section{Exemple motivant}
% ------------------------------------------------------------

Un psychologue compare les scores de bien-\^etre de deux groupes (m\'editation vs.\ contr\^ole) mais les donn\'ees sont ordinales (\'echelle de 1 \`a 10) et clairement non normales. Les tests $t$ et $F$ ne sont pas adapt\'es. Les m\'ethodes non param\'etriques offrent des alternatives valides sans hypoth\`ese de normalit\'e.

% ------------------------------------------------------------
\section{Pourquoi le non param\'etrique ?}
% ------------------------------------------------------------

\begin{itemize}
\item \textbf{Donn\'ees ordinales} ou de rang.
\item \textbf{Petits \'echantillons} o\`u la normalit\'e ne peut \^etre v\'erifi\'ee.
\item \textbf{Distributions fortement asym\'etriques} ou avec valeurs aberrantes.
\item \textbf{Distribution inconnue} : on ne veut pas imposer de mod\`ele param\'etrique.
\end{itemize}

\begin{attention}
Les tests non param\'etriques sont en g\'en\'eral \textbf{moins puissants} que leurs homologues param\'etriques quand les hypoth\`eses param\'etriques sont satisfaites. Mais ils sont plus \textbf{robustes}.
\end{attention}

% ------------------------------------------------------------
\section{Test des signes}
% ------------------------------------------------------------

\begin{definition}[Test des signes]
\index{test des signes}
Pour tester $H_0:\mathrm{m\'ediane}=m_0$, on compte $S^+=\#\{i:X_i>m_0\}$. Sous $H_0$ : $S^+\sim\mathcal{B}(n,1/2)$.
\end{definition}

% ------------------------------------------------------------
\section{Test de Wilcoxon pour un \'echantillon (rangs sign\'es)}
% ------------------------------------------------------------

\begin{definition}[Test de Wilcoxon sign\'e]
\index{Wilcoxon!rangs signes@rangs sign\'es}
\begin{enumerate}
\item Calculer $D_i=X_i-m_0$ et \'eliminer les $D_i=0$.
\item Ranger les $|D_i|$ par ordre croissant et affecter les rangs $R_i$.
\item $W^+=\sum_{\{i:D_i>0\}} R_i$.
\end{enumerate}
Sous $H_0:\mathrm{m\'ediane}=m_0$, $W^+$ a une distribution connue (tables ou approximation normale pour $n$ grand).
\end{definition}

% ------------------------------------------------------------
\section{Test de Mann--Whitney / Wilcoxon pour deux \'echantillons}
% ------------------------------------------------------------

\begin{definition}[Test de Mann--Whitney]
\index{Mann--Whitney}\index{Wilcoxon!deux echantillons@deux \'echantillons}
Pour comparer deux \'echantillons ind\'ependants $X_1,\ldots,X_{n_1}$ et $Y_1,\ldots,Y_{n_2}$ :
\begin{enumerate}
\item Combiner les deux \'echantillons et ranger toutes les observations.
\item $U_1=\sum_{i=1}^{n_1} R_i - \frac{n_1(n_1+1)}{2}$, o\`u $R_i$ est le rang de $X_i$.
\item Sous $H_0:F_X=F_Y$, $U$ a une distribution connue.
\end{enumerate}
\end{definition}

\begin{theoreme}[Approximation normale]
Pour $n_1,n_2$ grands :
\[
Z = \frac{U-\frac{n_1 n_2}{2}}{\sqrt{\frac{n_1 n_2(n_1+n_2+1)}{12}}} \xrightarrow{\mathcal{L}} \mathcal{N}(0,1).
\]
\end{theoreme}

% ------------------------------------------------------------
\section{Test de Kruskal--Wallis}
% ------------------------------------------------------------

\begin{definition}[Test de Kruskal--Wallis]
\index{Kruskal--Wallis}
G\'en\'eralisation du Mann--Whitney \`a $k$ groupes ($H_0$ : toutes les distributions sont identiques) :
\[
H = \frac{12}{n(n+1)}\sum_{j=1}^k \frac{R_j^2}{n_j} - 3(n+1),
\]
o\`u $R_j$ est la somme des rangs du groupe $j$. Sous $H_0$ : $H\sim\chi^2(k-1)$ asymptotiquement.
\end{definition}

% ------------------------------------------------------------
\section{Test de Kolmogorov--Smirnov}
% ------------------------------------------------------------

\begin{definition}[Test KS pour un \'echantillon]
\index{Kolmogorov--Smirnov}
Pour $H_0:F=F_0$ (loi sp\'ecifi\'ee) :
\[
D_n = \sup_{x}|\hat{F}_n(x)-F_0(x)|.
\]
La loi de $D_n$ sous $H_0$ est tabul\'ee (distribution de Kolmogorov).
\end{definition}

\begin{definition}[Test KS pour deux \'echantillons]
\[
D_{n_1,n_2} = \sup_x|\hat{F}_{n_1}(x)-\hat{G}_{n_2}(x)|.
\]
Teste si les deux \'echantillons proviennent de la m\^eme distribution.
\end{definition}

% ------------------------------------------------------------
\section{Test du $\chi^2$ d'ad\'equation}
% ------------------------------------------------------------

Rappel du th\'eor\`eme~\ref{thm:chi2-adequation} : ce test est d\'ej\`a non param\'etrique au sens o\`u il ne suppose pas de forme param\'etrique. Il compare les effectifs observ\'es aux effectifs attendus.

% ------------------------------------------------------------
\section{Estimation non param\'etrique de la densit\'e}
% ------------------------------------------------------------

\begin{definition}[Estimateur \`a noyau (KDE)]
\index{estimation a noyau@estimation \`a noyau}\index{KDE}
\[
\hat{f}_h(x) = \frac{1}{nh}\sum_{i=1}^n K\!\left(\frac{x-X_i}{h}\right),
\]
o\`u $K$ est un \textbf{noyau} (fonction int\'egrant \`a 1, typiquement gaussien) et $h>0$ est la \textbf{fen\^etre} (bandwidth).
\end{definition}

\begin{intuition}
$h$ petit $\to$ estimateur bruit\'e (faible biais, forte variance). $h$ grand $\to$ estimateur liss\'e (fort biais, faible variance). C'est le compromis biais-variance appliqu\'e \`a l'estimation de densit\'e.
\end{intuition}

% ------------------------------------------------------------
\section{Corr\'elation de Spearman}
% ------------------------------------------------------------

\begin{definition}[Coefficient de Spearman]
\index{Spearman}
Le coefficient de corr\'elation de Spearman $r_s$ est le coefficient de corr\'elation de Pearson appliqu\'e aux \textbf{rangs} :
\[
r_s = 1-\frac{6\sum d_i^2}{n(n^2-1)},
\]
o\`u $d_i=\mathrm{rang}(x_i)-\mathrm{rang}(y_i)$. Il mesure la monotonie de la relation (pas n\'ecessairement lin\'eaire).
\end{definition}

% ------------------------------------------------------------
\section{Impl\'ementation Python}
% ------------------------------------------------------------

\begin{python}
\begin{minted}{python}
import numpy as np
from scipy import stats
import matplotlib.pyplot as plt

np.random.seed(42)

# --- Test de Wilcoxon signe ---
avant = np.array([5.2, 4.8, 6.1, 5.5, 4.9, 5.8, 6.3, 5.1, 4.7, 5.4])
apres = np.array([4.8, 4.2, 5.5, 5.0, 4.3, 5.1, 5.7, 4.5, 4.1, 4.9])
stat_w, p_w = stats.wilcoxon(avant, apres, alternative='greater')
print(f"Wilcoxon signe: W={stat_w:.1f}, p={p_w:.4f}")

# --- Test de Mann-Whitney ---
groupe_A = np.array([23, 25, 28, 30, 27, 24, 26, 29, 31, 22])
groupe_B = np.array([18, 20, 22, 19, 21, 17, 23, 20, 19, 18])
u_stat, p_mw = stats.mannwhitneyu(groupe_A, groupe_B, alternative='two-sided')
print(f"Mann-Whitney: U={u_stat:.1f}, p={p_mw:.4f}")

# --- Test de Kruskal-Wallis ---
g1 = [6.4, 7.1, 6.8, 7.3, 6.9]
g2 = [8.2, 7.9, 8.5, 8.1, 8.4]
g3 = [5.5, 6.2, 5.8, 6.0, 5.7]
h_stat, p_kw = stats.kruskal(g1, g2, g3)
print(f"Kruskal-Wallis: H={h_stat:.4f}, p={p_kw:.4f}")

# --- Test de Kolmogorov-Smirnov ---
data = np.random.exponential(2, 50)
ks_stat, p_ks = stats.kstest(data, 'norm', args=(np.mean(data), np.std(data)))
print(f"KS (normalite): D={ks_stat:.4f}, p={p_ks:.4f}")

ks_stat2, p_ks2 = stats.kstest(data, 'expon', args=(0, np.mean(data)))
print(f"KS (exponentielle): D={ks_stat2:.4f}, p={p_ks2:.4f}")

# --- KDE ---
data_kde = np.concatenate([np.random.normal(-2, 0.8, 100),
                           np.random.normal(2, 1.2, 150)])
x_grid = np.linspace(-6, 7, 300)

fig, ax = plt.subplots(figsize=(10, 5))
ax.hist(data_kde, bins=30, density=True, alpha=0.3, label='Histogramme')
for bw in [0.2, 0.5, 1.0]:
    kde = stats.gaussian_kde(data_kde, bw_method=bw)
    ax.plot(x_grid, kde(x_grid), lw=2, label=f'KDE h={bw}')
ax.set_title("Estimation a noyau (KDE)")
ax.legend()
plt.savefig("ch11_kde.pdf")
plt.show()

# --- Correlation de Spearman ---
x = np.array([1, 2, 3, 4, 5, 6, 7, 8, 9, 10])
y = np.array([1, 4, 9, 16, 25, 36, 49, 64, 81, 100])  # y = x^2
r_pearson, _ = stats.pearsonr(x, y)
r_spearman, p_sp = stats.spearmanr(x, y)
print(f"Pearson: r={r_pearson:.4f}")
print(f"Spearman: r_s={r_spearman:.4f}, p={p_sp:.4f}")
\end{minted}
\end{python}

% ------------------------------------------------------------
\section{Exercices}
% ------------------------------------------------------------

\begin{exercice}[$\star$ -- Mann--Whitney]
Deux algorithmes sont test\'es sur 8 jeux de donn\'ees. Temps (en s) : A = [2.3, 3.1, 2.8, 3.5, 2.9, 3.2, 2.7, 3.0], B = [3.5, 4.2, 3.8, 4.0, 3.7, 4.1, 3.9, 3.6]. L'algorithme A est-il significativement plus rapide ?
\end{exercice}

\begin{exercice}[$\star\star$ -- Comparaison param\'etrique / non param\'etrique]
G\'en\'erer 1000 \'echantillons de taille $n=15$ d'une loi log-normale. Pour chaque \'echantillon, tester $H_0:\mathrm{m\'ediane}=1$ avec le test $t$ et le test de Wilcoxon. Comparer les taux d'erreur de type I.
\end{exercice}

\begin{exercice}[$\star\star$ -- KDE]
Impl\'ementer un estimateur \`a noyau gaussien. Tester sur un m\'elange de deux gaussiennes et \'etudier l'effet de $h$ sur la qualit\'e de l'estimation (MISE).
\end{exercice}

\begin{exercice}[$\star\star\star$ -- Projet]
Collecter des donn\'ees r\'eelles non normales (par ex. revenus, dur\'ees d'attente). Appliquer les m\'ethodes non param\'etriques vues dans ce chapitre. Comparer avec les r\'esultats des tests param\'etriques et discuter.
\end{exercice}

% ------------------------------------------------------------
\begin{formules}
\begin{align*}
U &= \sum R_i - \frac{n_1(n_1+1)}{2} & H &= \frac{12}{n(n+1)}\sum\frac{R_j^2}{n_j}-3(n+1) \\
D_n &= \sup_x|\hat{F}_n(x)-F_0(x)| & r_s &= 1-\frac{6\sum d_i^2}{n(n^2-1)} \\
\hat{f}_h(x) &= \frac{1}{nh}\sum K\!\left(\frac{x-X_i}{h}\right)
\end{align*}
\end{formules}
